\textcolor{red}{(Este capítulo em geral não deve ultrapassar o total de 3 páginas.)}

\section{Contribuições}

\textcolor{red}{Nesta seção, o aluno deve destacar as principais contribuições do trabalho realizado, tanto no nível profissional quanto pessoal. Descreva conclusões a respeito do trabalho desenvolvido e experiências adquiridas.}

\section{Considerações Sobre o Curso de Graduação}

\textcolor{red}{Pode começar descrevendo como as disciplinas do curso ou atividades desenvolvidas durante o curso de graduação se relacionam ou contribuíram para o projeto desenvolvido. 
 Depois, coloque aqui as suas críticas ao curso de graduação, pontuando os pontos fracos, (e porque não) os fortes também, que limitaram (ou contribuíram para) suas ações no desenvolvimento do trabalho. Faça sugestões que você achar pertinente para a melhora do curso, disciplinas ou temas novos a serem abordados, etc.  }

\section{Trabalhos Futuros}

\textcolor{red}{Nesta seção, o aluno pode apresentar atividades e sugestões a serem realizadas para a continuação do trabalho realizado neste projeto. Esta seção é opcional.}

\section{Referências (Seção apenas explicativa, não deve ir no trabalho)}

\textcolor{red}{Observação 1: As referências devem estar em ordem alfabética pelo sobrenome do primeiro autor. Observação 2: Todos os documentos referenciados nesta seção (livros, artigos, relatórios técnicos, sites, entre outros) deverão ter sido citados no texto. Observação 3: Todos as citações no decorrer do texto deverão ser listadas nesta seção. Observação 4: Todos as referências devem ser identificáveis quanto a forma de publicação: por exemplo, se for um livro, deve ter a Editora, se for um artigo, deve ter o nome do periódico, se for uma tese ou dissertação, isso dever estar claro e deve ter o nome da Universidade onde foi defendida. Veja a observação 5 a respeito disso, pois seguindo a norma não tem como errar. Observação 5: É obrigatório que a lista de referências esteja de acordo com a norma NBR-6023/2002 para referenciação bibliográfica.}